\chapter{Servizi SOAP/WSDL}

\section{Introduzione}

I servizi RPC (Remote Procedure Call) basati su SOAP/WSDL rappresentano un approccio più formale e strutturato alla definizione e implementazione di web service rispetto ai servizi RESTful. Si basano su un insieme di tecnologie standardizzate che permettono l'interoperabilità tra sistemi differenti.

Le principali tecnologie usate sono:

\begin{itemize}
    \item \textbf{XML Schema}: per definire la struttura e i tipi di dato dei documenti XML scambiati.
    \item \textbf{UDDI} (Universal Description, Discovery and Integration): per la pubblicazione e ricerca di web service.
    \item \textbf{WSDL} (Web Services Description Language): per descrivere l'interfaccia del web service.
    \item \textbf{SOAP} (Simple Object Access Protocol): il protocollo per lo scambio di messaggi.
    \item \textbf{HTTP}: come protocollo di trasporto.
\end{itemize}

\section{XML Schema}

\subsection{Cosa permette di fare XML Schema}

XML Schema permette di:

\begin{itemize}
    \item Definire gli elementi (tag) che possono apparire in un documento.
    \item Definire gli attributi che possono apparire in un elemento.
    \item Definire quali elementi devono essere inseriti in altri elementi (chiamati \emph{child}).
    \item Definire il numero degli elementi child.
    \item Definire quando un elemento deve essere vuoto o può contenere testo, elementi, oppure entrambi.
    \item Definire il tipo per ogni elemento e per gli attributi (intero, stringa, ecc., ma anche tipi personalizzati).
    \item Definire i valori di default o fissi per elementi ed attributi.
\end{itemize}

\subsection{Esempio di documento XML}

\begin{lstlisting}[language=XML]
<?xml version="1.0"?>
<libro>
    <titolo>Al di la' del bene e del male</titolo>
    <autore>Nietzsche</autore>
    <numeroDiPagine>272</numeroDiPagine>
</libro>
\end{lstlisting}

\subsection{Esempio di XML Schema}

Il corrispondente XML Schema che definisce la struttura del documento precedente:

\begin{lstlisting}[language=XML]
<?xml version="1.0"?>
<!-- iniziamo lo schema -->
<xs:schema xmlns:xs="http://www.w3.org/2001/XMLSchema"
           targetNamespace="localhost/wsdl"
           xmlns="localhost/wsdle"
           elementFormDefault="qualified">
    <!-- definiamo l'elemento libro -->
    <xs:element name="libro">
        <xs:complexType>
            <xs:all>
                <!-- definiamo i vari elementi child di libro -->
                <xs:element name="titolo" type="xs:string" />
                <xs:element name="autore" type="xs:string" />
                <xs:element name="numeroDiPagine" type="xs:integer" />
            </xs:all>
        </xs:complexType>
    </xs:element>
</xs:schema>
\end{lstlisting}

\subsection{Spiegazione degli attributi principali}

\begin{itemize}
    \item \texttt{xmlns:xs="..."}: stabilisce che gli elementi che iniziano con \texttt{xs} provengono dal namespace \url{http://www.w3.org/2001/XMLSchema}.
    \item \texttt{targetNamespace="..."}: definisce a quale namespace appartengono gli elementi definiti con questo schema.
    \item \texttt{xmlns="..."}: definisce il namespace di default (quello applicato se non ne esplicitiamo uno).
    \item \texttt{elementFormDefault="qualified"}: indica che ogni elemento usato nel documento XML che utilizzerà questo schema dovrà essere qualificato.
\end{itemize}

\section{Elementi di un documento WSDL}

Un documento WSDL ha la seguente struttura generale:

\begin{lstlisting}[language=XML]
<definitions>
    <types>
        <!-- definizione dei tipi di dato utilizzati -->
    </types>
    <message>
        <!-- definizione di uno dei messaggi impiegati dal
             web service per comunicare con l'applicazione client -->
    </message>
    <!-- naturalmente puo' esistere piu' di un elemento message
         all'interno del documento -->
    <portType>
        <!-- definisce una "porta" e le operazioni che possono
             essere eseguite dal web service.
             Definisce inoltre i messaggi coinvolti nelle
             operazioni elencate -->
    </portType>
    <binding>
        <!-- definisce il formato del messaggio ed i
             dettagli di protocollo per ogni porta -->
    </binding>
</definitions>
\end{lstlisting}

Analizziamo nel dettaglio ciascuno degli elementi.

\subsection{Types (Tipi)}

La sezione \texttt{<types>} definisce i tipi di dato utilizzati nei messaggi del web service, generalmente tramite uno XML Schema:

\begin{lstlisting}[language=XML]
<types>
    <xs:schema targetNamespace="http://localhost/wsdl">
        <xs:element name="email" type="xs:string"/>
        <xs:element name="password" type="xs:string"/>
        <xs:element name="respText" type="xs:string"/>
        <xs:element name="utente">
            <xs:complexType>
                <xs:all>
                    <xs:element name="email" type="xs:string"/>
                    <xs:element name="password" type="xs:string"/>
                </xs:all>
            </xs:complexType>
        </xs:element>
    </xs:schema>
</types>
\end{lstlisting}

\subsection{Messages (Messaggi)}

La sezione \texttt{<message>} definisce i messaggi che il web service scambia con il client. Ogni operazione ha tipicamente un messaggio di \emph{request} (inviato al server) e uno di \emph{response} (restituito dal server):

\begin{lstlisting}[language=XML]
<message name="get_userRequest">
    <part name="email" type="xs:string" />
</message>
<message name="get_userResponse">
    <part name="utente" type="xs:utente" />
</message>
<message name="save_userRequest">
    <part name="email" type="xs:string" />
    <part name="password" type="xs:string" />
</message>
<message name="save_userResponse">
    <part name="respText" type="xs:string" />
</message>
\end{lstlisting}

\subsection{PortType}

L'elemento \texttt{<portType>} definisce le operazioni che il web service può eseguire e i messaggi coinvolti in ciascuna operazione:

\begin{lstlisting}[language=XML]
<portType name="gestioneUtentiType">
    <operation name="getUserByEmail">
        <input message="tns:get_userRequest" />
        <output message="tns:get_userResponse" />
    </operation>
    <operation name="saveUser">
        <input message="tns:save_userRequest" />
        <output message="tns:save_userResponse" />
    </operation>
</portType>
\end{lstlisting}

Ogni \texttt{<operation>} indica un messaggio di input (la richiesta) e uno di output (la risposta).

\subsection{Binding}

L'elemento \texttt{<binding>} definisce il formato del messaggio e i dettagli di protocollo per ogni porta. Collega il \texttt{portType} a un protocollo di trasporto concreto (tipicamente SOAP su HTTP):

\begin{lstlisting}[language=XML]
<binding name="gestioneUtentiBinding" type="tns:gestioneUtentiType">
    <soap:binding style="rpc"
        transport="http://schemas.xmlsoap.org/soap/http"/>
    <operation name="getUserByEmail">
        <soap:operation
            soapAction="localhost/wsdl/service.php/getUserByEmail"
            style="rpc"/>
        <input>
            <soap:body use="encoded"
                namespace="localhost/wsdl"
                encodingStyle="http://schemas.xmlsoap.org/soap/encoding/"/>
        </input>
        <output>
            <soap:body use="encoded"
                namespace="localhost/wsdl"
                encodingStyle="http://schemas.xmlsoap.org/soap/encoding/"/>
        </output>
    </operation>
    <operation name="saveUser">
        <soap:operation
            soapAction="localhost/wsdl/service.php/saveUser"
            style="rpc"/>
        <input>
            <soap:body use="encoded"
                namespace="localhost/wsdl"
                encodingStyle="http://schemas.xmlsoap.org/soap/encoding/"/>
        </input>
        <output>
            <soap:body use="encoded"
                namespace="localhost/wsdl"
                encodingStyle="http://schemas.xmlsoap.org/soap/encoding/"/>
        </output>
    </operation>
</binding>
\end{lstlisting}

\subsection{Service}

L'elemento \texttt{<service>} definisce l'endpoint del web service, cioè l'indirizzo al quale il servizio è raggiungibile:

\begin{lstlisting}[language=XML]
<service name="ServizioUtenti">
    <port name="GestioneUtenti"
          binding="tns:gestioneUtentiBinding">
        <soap:address location="localhost/wsdl/service.php"/>
    </port>
</service>
\end{lstlisting}

\section{Considerazioni finali}

Un documento WSDL completo fornisce tutte le informazioni necessarie perché un client possa invocare il web service:

\begin{itemize}
    \item Quali operazioni sono disponibili (\texttt{portType}).
    \item Quali messaggi vengono scambiati in ogni operazione (\texttt{message}).
    \item Come sono strutturati i dati nei messaggi (\texttt{types}).
    \item Come devono essere trasmessi i messaggi sulla rete (\texttt{binding}).
    \item Dove (a quale URL) si trova il servizio (\texttt{service}).
\end{itemize}

Rispetto ai servizi RESTful, i servizi SOAP/WSDL offrono maggiore formalità e rigore nella specifica dell'interfaccia, ma sono generalmente più complessi da progettare, implementare e utilizzare.
